%?deps @bib-commons

\documentclass{beamer}
\usepackage{amsmath,amssymb,amsthm,hyperref,mathrsfs,mathtools,mleftright,xcolor}
\usepackage[en-GB]{datetime2}
\usepackage{bib-commons}

\usetheme[shadow=false]{Boadilla}

\setbeamerfont{institute}{size=\small}
\addtobeamertemplate{institute}{\vskip-1.1em}{}

\makeatletter
% scale footline sections
\setbeamertemplate{footline}{
  \leavevmode%
  \hbox{%
    \begin{beamercolorbox}[wd=.25\paperwidth,ht=2.25ex,dp=1ex,center]{author in head/foot}%
      \usebeamerfont{author in head/foot}\insertshortauthor%
    \end{beamercolorbox}%
    \begin{beamercolorbox}[wd=.5\paperwidth,ht=2.25ex,dp=1ex,center]{title in head/foot}%
      \usebeamerfont{title in head/foot}\insertshorttitle%
    \end{beamercolorbox}%
    \begin{beamercolorbox}[wd=.25\paperwidth,ht=2.25ex,dp=1ex,leftskip=2ex,rightskip=2ex,sep=0pt]{date in head/foot}%
      \hfill%
      \usebeamerfont{date in head/foot}%
      \insertshortdate{}%
      \hfill%
      \usebeamercolor[fg]{page number in head/foot}%
      \usebeamerfont{page number in head/foot}%
      \usebeamertemplate{page number in head/foot}%
    \end{beamercolorbox}}%
  \vskip0pt}

\newtheorem{named@thm}{\curr@thmname}
\newenvironment{namedthm}[1]
  {\def\curr@thmname{#1}\begin{named@thm}}
  {\end{named@thm}}
\makeatother

\newcommand{\tuple}[1]{\left\langle #1 \right\rangle}

\newcommand{\ZF}{\mathrm{ZF}}
\newcommand{\IZF}{\mathrm{IZF}}
\newcommand{\ZFC}{\mathrm{ZFC}}
\newcommand{\Tr}{\mathrm{Tr}}
\newcommand{\Ord}{\mathrm{Ord}}

\DeclareMathOperator{\Def}{Def}

\newcommand{\exampletext}[1]{\textcolor{green!50!black}{#1}}

\makeatletter
\newsavebox\myin@eqfst
\newsavebox\myin@eqsnd
\def\myineq{\mathrel{\mathpalette\myin@eq\relax}}
\def\myin@eq#1#2{%
  \setbox\myin@eqfst=\hbox{\m@th$#1\in$}%
  \setbox\myin@eqsnd=\hb@xt@\wd\myin@eqfst{\m@th$#1\dabar@$\hss$#1\dabar@$}%
  \vcenter{\offinterlineskip%
    \box\myin@eqfst\vskip-.1\ht\myin@eqsnd%
    \copy\myin@eqsnd\vskip-.5\ht\myin@eqsnd}}
\makeatother

\title[Indexed hierarchies and the ordinal class]{Indexed hierarchies and internal structure in the\\ordinal class in intuitionistic set theories}
\author{Shuwei Wang}
\institute{University of Leeds}
\date[LC26]{\small Logic Colloquium 2026\\\DTMdate{2026-06-30}}

\begin{document}

\begin{frame}
  \titlepage
\end{frame}

\begin{frame}{Ordinals and ordinal-indexed hierarchies}
  In $\ZF$, we say
  \[\Ord = \left\{\alpha : \text{$\alpha$ is transitive} \land \forall \beta \in \alpha \ \text{$\beta$ is transitive}\right\}.\]

  \pause

  Fix a set-like partial order $\mathscr{R}$ on the class $\Tr$ of all transitive sets, which is no coarser than $\subseteq$ yet no finer than $\myineq$, i.e.
  \[x = y \lor x \in y \Longrightarrow x \mathscr{R} y; \quad x \mathscr{R} y \Longrightarrow x \subseteq y,\]
  and we consider the \emph{$\mathscr{R}$-predecessor} set
  \[P_\mathscr{R}\mleft(x\mright) = \left\{y : y \mathop{\mathscr{R}} x\right\}.\]
  E.g.\ when $\mathscr{R} = {\subseteq}$, $P_\mathscr{R}\mleft(x\mright) = \mathcal{P}\mleft(x\mright)$; when $\mathscr{R} = {\myineq}$, $P_\mathscr{R}\mleft(x\mright) = x^+ = x \cup \left\{x\right\}$. We can also get the first-order definable subsets $\Def\mleft(x\mright)$, etc.
\end{frame}

\begin{frame}{Ordinals and ordinal-indexed hierarchies}
  Fix a set-like partial order $\mathscr{R}$ on the class $\Tr$ of all transitive sets, which is no coarser than $\subseteq$ yet no finer than $\myineq$. We can define
  \begin{align*}
    H_\mathscr{R}\mleft(\varnothing\mright) & = \varnothing,                                                                                      \\
    H_\mathscr{R}\mleft(\alpha^+\mright)    & = P_\mathscr{R}\mleft(H_\mathscr{R}\mleft(\alpha\mright)\mright),                                   \\
    H_\mathscr{R}\mleft(\lambda\mright)     & = \bigcup_{\alpha \in \lambda} H_\mathscr{R}\mleft(\alpha\mright) \quad \text{for limit $\lambda$}.
  \end{align*}
  E.g.\ $V_\alpha$, $L_\alpha$, etc.

  \pause

  \vspace{1em}

  Alternatively, for any set $x$, we can define
  \[H_\mathscr{R}\mleft(x\mright) = \bigcup_{y \in x} P_\mathscr{R}\mleft(H_\mathscr{R}\mleft(y\mright)\mright).\]
  One easily verifies that this gives the same results on ordinals.
\end{frame}

\begin{frame}{Ordinals and ordinal-indexed hierarchies}
  Fix a set-like partial order $\mathscr{R}$ on the class $\Tr$ of all transitive sets, which is no coarser than $\subseteq$ yet no finer than $\myineq$.

  \begin{block}{Question}
    Is there a maximal class $O_\mathscr{R} \subseteq V$ on which $H_\mathscr{R}$ is an injection?
  \end{block}

  `Maximal' means for any $x \in V$, there exists $y \in O_\mathscr{R}$ such that $H_\mathscr{R}\mleft(x\mright) = H_\mathscr{R}\mleft(y\mright)$.

  \pause

  \vspace{1em}

  \exampletext{Answer ($\ZF$).} Always take $O_\mathscr{R} = \Ord$.

  \vspace{0.4em}

  Injectivity follows from linearity of ordinals (note $x \in y$ implies $H_\mathscr{R}\mleft(x\mright) \in H_\mathscr{R}\mleft(x^+\mright) \subseteq H_\mathscr{R}\mleft(y\mright)$).

  \vspace{0.4em}

  For any $x \in V$, simply observe $H_\mathscr{R}\mleft(x\mright) = H_\mathscr{R}\mleft(\mathrm{rank}\mleft(x\mright)\mright)$.
\end{frame}

\begin{frame}{Intuitionistic set theory}
  Intuitionistic Zermelo--Fraenke set theory $\IZF$ is:

  \vspace{0.4em}

  $\ZF$, minus Excluded Middle, with Axiom of Foundation/Regularity replaced by Set Induction:
  \[\forall x \left(\forall y \in x \ \varphi\mleft(y\mright) \rightarrow \varphi\mleft(x\mright)\right) \rightarrow \forall x \ \varphi\mleft(x\mright).\]

  \pause

  \begin{namedthm}{Wellknown folklore}
    $\IZF$ is consistent with $\Ord$ being non-linear, or formally
    \[\exists \alpha, \beta \in \Ord \left(\beta \subseteq \alpha \land \beta \not\in \alpha \land \beta \neq \alpha\right).\]
  \end{namedthm}

  \vspace{0.4em}

  This is easily implementable in multiple types of `intuitionistic models', e.g.\ models on Kripke frames, realisability models or Heyting-valued models.
\end{frame}

\begin{frame}{Intuitionistic set theory}
  \begin{corollary}
    $\IZF$ is consistent with the class function $\alpha \mapsto V_\alpha$ being non-injective.
  \end{corollary}

  \vspace{0.4em}

  \exampletext{Proof.} Take $\alpha, \beta \in \Ord$ such that $\beta \subseteq \alpha$, $\beta \not\in \alpha$ and $\beta \neq \alpha$. Consider
  \[V_{\alpha^+} = \mathcal{P}\mleft(V_\alpha\mright).\]
  But $V_\beta \subseteq V_\alpha$, i.e.\ $\mathcal{P}\mleft(V_\beta\mright) \subseteq \mathcal{P}\mleft(V_\alpha\mright)$, thus
  \[V_{\alpha^+ \cup \beta^+} = \mathcal{P}\mleft(V_\alpha\mright) \cup \mathcal{P}\mleft(V_\beta\mright) = V_{\alpha^+}.\]
  Clearly $\beta \not\in \alpha^+$, so $\alpha^+ \neq \alpha^+ \cup \beta^+$. \qed
\end{frame}

\begin{frame}{The intuition}
  \begin{block}{Question}
    Is there a maximal class $O_\mathscr{R} \subseteq V$ on which $H_\mathscr{R}$ is an injection?
  \end{block}

  \vspace{0.4em}

  It is still reasonable to expect $O_\mathscr{R} \subseteq \Ord$. Ordinals in $O_\mathscr{R}$ should satisfy some additional `closure properties'.

  \pause

  \vspace{1em}

  \exampletext{Example (Paul Taylor, 1996\nocite{taylor96-intuitionistic-ordinals})}: We say an ordinal $\alpha \in \Ord$ is \emph{plump} if
  \begin{itemize}
    \item every $\beta \in \alpha$ is plump;
    \item for every $\gamma \in \Ord$, if $\gamma$ is plump and $\gamma \subseteq \beta$ for some $\beta \in \alpha$, then $\gamma \in \alpha$.
  \end{itemize}

  As it turns out, the class of plump ordinals is precisely a maximal class on which $\alpha \mapsto V_\alpha$ is injective.
\end{frame}

\begin{frame}{The intuition}
  A similar idea comes from Robert Lubarsky, \emph{Intuitionistic L} (1993)\nocite{lubarsky93-intuitionistic-l}:
  \begin{theorem}
    Although it is open whether $\IZF \vdash \Ord \subseteq L$, we can avoid it and still prove $\IZF \vdash \left(V = L\right)^L$.
  \end{theorem}

  \vspace{0.4em}

  \exampletext{Proof.} For any $\alpha \in \Ord$, we construct $\alpha^* \in L$ as
  \[\alpha^* = \left\{\beta \in \Ord : L_{\left(\alpha_\text{aug} +_\mathrm{H} n\right)^-} \vDash \mathrm{def}\mleft(L_\beta\mright) \subseteq L_\alpha\right\},\]
  such that $L_\alpha = L_{\alpha^*} \in L$. \qed

  \vspace{1em}

  Note that $\mathrm{def}\mleft(L_\beta\mright) \subseteq L_\alpha$ can be equivalently written as $L_\beta \in L_\alpha$, since if $L_\beta \in \mathrm{def}\mleft(L_\gamma\mright)$ for some $\gamma \in \alpha$, then also
  \[\mathrm{def}\mleft(L_\beta\mright) \subseteq \mathrm{def}\mleft(L_\gamma\mright) \subseteq L_\alpha.\]
\end{frame}

\begin{frame}{The formal construction}
  We now directly define
  \begin{align*}
    O_\mathscr{R} = \{\alpha \in \Ord : {} & \forall \beta \in \alpha^+ \ \forall \gamma \in \Ord                                                                                                                                                      \\
                                           & \left(H_\mathscr{R}\mleft(\gamma\mright) \in H_\mathscr{R}\mleft(\beta\mright) \rightarrow \exists \gamma' \in \beta \ H_\mathscr{R}\mleft(\gamma'\mright) = H_\mathscr{R}\mleft(\gamma\mright)\right)\}.
  \end{align*}

  \pause

  \begin{namedthm}{Proposition}
    $H_\mathscr{R}$ is injective on $O_\mathscr{R}$.
  \end{namedthm}

  \vspace{0.4em}

  \exampletext{Proof.} Use induction. Suppose $\beta, \beta' \in O_\mathscr{R}$ and $H_\mathscr{R}\mleft(\beta\mright) = H_\mathscr{R}\mleft(\beta'\mright)$, we want to show $\beta = \beta'$. For $\gamma \in \beta$,
  \[H_\mathscr{R}\mleft(\gamma\mright) \in H_\mathscr{R}\mleft(\beta\mright) = H_\mathscr{R}\mleft(\beta'\mright),\]
  i.e.\ $H_\mathscr{R}\mleft(\gamma\mright) = H_\mathscr{R}\mleft(\gamma'\mright)$ for $\gamma' \in \beta'$. By inductive hypothesis, $\gamma = \gamma' \in \beta'$. The converse is symmetrical. \qed
\end{frame}

\begin{frame}{$\mathscr{R}$-successors}
  Are there many elements in $O_\mathscr{R}$?

  \vspace{0.4em}

  We want to first find the $\mathscr{R}$-successor of $\alpha \in O_\mathscr{R}$, i.e.\ some $\alpha^{\mathscr{R}+} \in O_\mathscr{R}$ such that
  \[H_\mathscr{R}\mleft(\alpha^{\mathscr{R}+}\mright) = H_\mathscr{R}\mleft(\alpha^+\mright) = P_\mathscr{R}\mleft(H_\mathscr{R}\mleft(\alpha\mright)\mright).\]

  \pause

  \vspace{0.4em}

  We consider the set
  \[s = \left\{\left\{\gamma \in \alpha : H_\mathscr{R}\mleft(\gamma\mright) \in x\right\} : x \in P_\mathscr{R}\mleft(H_\mathscr{R}\mleft(\alpha\mright)\mright)\right\}.\]
  For any $\beta \in \Ord$, if $H_\mathscr{R}\mleft(\beta\mright) \in P_\mathscr{R}\mleft(H_\mathscr{R}\mleft(\alpha\mright)\mright)$, then for each $\gamma \in \beta$, there is $\gamma' \in \alpha$ such that $H_\mathscr{R}\mleft(\gamma'\mright) = H_\mathscr{R}\mleft(\gamma\mright)$. Thus
  \[H_\mathscr{R}\mleft(\beta\mright) = \bigcup_{\gamma \in \beta} P_\mathscr{R}\mleft(H_\mathscr{R}\mleft(\gamma\mright)\mright) = \bigcup_{\gamma' \in t} P_\mathscr{R}\mleft(H_\mathscr{R}\mleft(\gamma'\mright)\mright) = H_\mathscr{R}\mleft(t\mright)\]
  for $t = \left\{\gamma \in \alpha : H_\mathscr{R}\mleft(\gamma\mright) \in H_\mathscr{R}\mleft(\beta\mright)\right\} \in s$. Easy to verify $t \in O_\mathscr{R}$.
\end{frame}

\begin{frame}{$\mathscr{R}$-successors}
  For $\alpha \in O_\mathscr{R}$, consider the set
  \[s = \left\{\left\{\gamma \in \alpha : H_\mathscr{R}\mleft(\gamma\mright) \in x\right\} : x \in P_\mathscr{R}\mleft(H_\mathscr{R}\mleft(\alpha\mright)\mright)\right\},\]
  it suffices to define $\alpha^{\mathscr{R}+} = \left\{\beta \in s : \beta \in O_\mathscr{R}, H_\mathscr{R}\mleft(\beta\mright) \mathop{\mathscr{R}} H_\mathscr{R}\mleft(\alpha\mright)\right\}$.

  \begin{lemma}
    For any $\alpha \in O_\mathscr{R}$, $\alpha^{\mathscr{R}+} \in O_\mathscr{R}$ is the least element in $O_\mathscr{R}$ that contains $\alpha$.
  \end{lemma}

  \pause

  \begin{lemma}
    For $s \subseteqq O_\mathscr{R}$, $\sup s = \bigcup s \in O_\mathscr{R}$.
  \end{lemma}

  \vspace{0.4em}

  So we can define $\mathrm{rk}_\mathscr{R} : V \rightarrow O_\mathscr{R}$ as
  \[\mathrm{rk}_\mathscr{R}\mleft(x\mright) = \bigcup_{y \in x} \mathrm{rk}_\mathscr{R}\mleft(y\mright)^{\mathscr{R}+}.\]
\end{frame}

\begin{frame}{Maximality of $O_\mathscr{R}$}
  \[\mathrm{rk}_\mathscr{R}\mleft(x\mright) = \bigcup_{y \in x} \mathrm{rk}_\mathscr{R}\mleft(y\mright)^{\mathscr{R}+}.\]

  Using $H_\mathscr{R}\mleft(\alpha^{\mathscr{R}+}\mright) = P_\mathscr{R}\mleft(H_\mathscr{R}\mleft(\alpha\mright)\mright)$, we prove by induction that:
  \begin{namedthm}{Proposition}
    For any set $x$,
    \[H_\mathscr{R}\mleft(x\mright) = H_\mathscr{R}\mleft(\mathrm{rk}_\mathscr{R}\mleft(x\mright)\mright).\]
  \end{namedthm}

  \begin{corollary}
    For any set $x$, there exists $\alpha \in O_\mathscr{R}$ such that $H_\mathscr{R}\mleft(x\mright) = H_\mathscr{R}\mleft(\alpha\mright)$, i.e.\ $O_\mathscr{R}$ is a maximal class on which $H_\mathscr{R}$ is injective.
  \end{corollary}
\end{frame}

\begin{frame}{Consistency of class separations}
  \begin{itemize}
    \item $\ZF$ proves that $O_\mathscr{R} = \Ord$ for all $\mathscr{R}$.

    \item When $\mathscr{R} = {\myineq}$, $H_\mathscr{R}\mleft(x\mright) = \mathrm{rk}_\Ord\mleft(x\mright)$, so we always have $O_{\myineq} = \Ord$.

          \pause

    \item We showed at the beginning that $\IZF$ is consistent with $\alpha \neq \beta \land V_\alpha = V_\beta$, so it is consistent with $O_\subseteq \neq \Ord = O_{\myineq}$.

          \pause

    \item Let $\mathscr{R}_\mathrm{def}$ be (first-order) definable subset relation, so that $H_{\mathscr{R}_\mathrm{def}}\mleft(\alpha\mright) = L_\alpha$. It is consistent that $2 \not\in O_\subseteq$ (when $1$ has undecidable subsets, i.e.\ $\alpha \subseteq 1$ with $\varnothing \in \alpha \lor \varnothing \not\in \alpha$ failing to hold), but $2 = L_2 \in O_{\mathscr{R}_\mathrm{def}}$ always, so consistently $O_\subseteq \neq O_{\mathscr{R}_\mathrm{def}}$.

          \pause

    \item Using realisability arguments, we can verify
          \begin{theorem}[W., upcoming PhD thesis]
            $\IZF$ is consistent with $\alpha \neq \beta \land L_\alpha = L_\beta$.
          \end{theorem}
          Thus also consistently $O_{\mathscr{R}_\mathrm{def}} \neq \Ord = O_{\myineq}$.
  \end{itemize}
\end{frame}

\begin{frame}{Plump ordinals are useful!}
  It remains open from Lubarsky (1993) whether $\IZF \vdash \Ord \subseteq L$. However,
  \begin{theorem}[W., 2026\nocite{wang-sw26-plump}]
    $O_\subseteq \subseteq L$. More specifically, for any $\alpha \in O_\subseteq$,
    \[\alpha = \left\{\beta \in L_\alpha : L_\alpha \vDash \beta \in O_\subseteq\right\}.\]
  \end{theorem}
  (We need to phrase $O_\subseteq$ as Taylor's plump ordinals, to avoid mentioning the objects $H_\subseteq\mleft(\beta\mright) = V_\beta$ in $L_\alpha$.)

  \pause

  \vspace{0.4em}

  Combined with the use of incomparable ordinals, it follows that
  \begin{theorem}[W., 2026]
    Fix any set $x$ in any $\ZFC$ universe, then there is a Heyting algebra $\mathbb{H}$ such that the Heyting extension satisfies
    \[V^\mathbb{H} \vDash \mathcal{P}\mleft(\check{x}\mright) \in L.\]
  \end{theorem}
\end{frame}

\begin{frame}{Ordinal arithmetic}
  Fix any $\mathscr{R}$, we can let
  \[\alpha +_\mathscr{R} \beta = \alpha \cup \bigcup_{\gamma \in \beta} \left(\alpha +_\mathscr{R} \gamma\right)^{\mathscr{R}+}, \quad \alpha \cdot_\mathscr{R} \beta = \bigcup_{\gamma \in \beta} \left(\alpha \cdot_\mathscr{R} \gamma +_\mathscr{R} \alpha\right).\]

  \pause

  \begin{lemma}
    For any $\alpha, \beta, \gamma \in \Ord = O_{\myineq}$, if $\alpha + \beta = \alpha + \gamma$, then $\beta = \gamma$; if $\alpha + \beta \in \alpha + \gamma$, then $\beta \in \gamma$.
  \end{lemma}

  \begin{lemma}[W., 2026]
    For any $\alpha, \beta, \gamma \in O_\subseteq$, if $\alpha +_\subseteq \beta = \alpha +_\subseteq \gamma$, then $\beta = \gamma$; if $\alpha +_\subseteq \beta \subseteq \alpha +_\subseteq \gamma$, then $\beta \subseteq \gamma$.
  \end{lemma}

  \pause

  \vspace{0.4em}

  These proofs are different (and indeed, the $O_\subseteq$ one is even simpler), and I do not know any uniform way to do this for arbitrary $\mathscr{R}$.
\end{frame}

\begin{frame}[b]
  \relax % otherwise \begin{frame} take the next line as title
  {\Large\usebeamercolor[fg]{frametitle} Thank you!}

  \vfill

  \let\oldbiblist\biblist
  \def\biblist{\oldbiblist\footnotesize}
  \hrule width 0.333\textwidth
  \bibmain
\end{frame}

\end{document}